\documentclass[../DoAn.tex]{subfiles}
\begin{document}

\section{Kết luận}
\label{section:6.1}

Sau quá trình nghiên cứu, phân tích và xây dựng, đồ án đã hoàn thành phần mềm Quản lý Khen thưởng tại Học viện Khoa học Quân sự (PM QLKT) với các kết quả cụ thể như sau.

\textbf{Về mặt nghiệp vụ}, hệ thống bao quát toàn bộ bảy nhóm khen thưởng đặc thù của môi trường quân sự: khen thưởng cá nhân hằng năm theo chuỗi CSTĐCS -- BKBTBQP -- CSTĐTQ -- BKTTCP, khen thưởng đơn vị hằng năm theo chuỗi ĐVQT -- BKBTBQP -- BKTTCP, Huy chương Chiến sĩ vẻ vang (HCCSVV) theo các mốc 10/15/20 năm, Huân chương Bảo vệ Tổ quốc (HCBVTQ) dựa trên thời gian giữ chức vụ theo nhóm hệ số, Kỷ niệm chương vì sự nghiệp xây dựng QĐNDVN, Huy chương Quân kỳ quyết thắng (HCQKQT, 25 năm phục vụ) và thành tích nghiên cứu khoa học. Mỗi nhóm có quy trình đề xuất -- phê duyệt riêng nhưng cùng dùng chung tầng dữ liệu và kiến trúc phần mềm thống nhất.

\textbf{Về mặt kỹ thuật}, hệ thống được hiện thực hoá theo kiến trúc phân tầng sáu lớp (Route $\rightarrow$ Middleware $\rightarrow$ Controller $\rightarrow$ Service $\rightarrow$ Repository $\rightarrow$ Prisma), với cơ sở dữ liệu PostgreSQL có 23 model chuẩn hoá, lớp xác thực JWT có cơ chế làm mới luân phiên, hệ thống ghi nhật ký kiểm toán đầy đủ qua middleware \texttt{auditLog}. Kho mã nguồn đi kèm 75 bộ kiểm thử với 890 ca kiểm thử đơn vị đạt 100\,\% thành công, phủ toàn bộ rule chuỗi cá nhân -- đơn vị, các kịch bản lỡ đợt, cửa sổ trượt và lifetime block.

\textbf{Về đóng góp nổi bật}, đồ án đã trừu tượng hoá rule chuỗi danh hiệu thành đối tượng cấu hình \texttt{ChainAwardConfig} (số năm chu kỳ, các cờ tiền điều kiện, có yêu cầu NCKH hay không, lifetime hay non-lifetime) và đặt vào hai mảng \texttt{PERSONAL\_CHAIN\_AWARDS} / \texttt{UNIT\_CHAIN\_AWARDS}. Hàm xét duy nhất \texttt{checkChainEligibility} đọc cấu hình này --- thiết kế giúp việc thêm tier mới chỉ cần bổ sung phần tử vào mảng mà không phải sửa logic. Tương tự, bảy loại đề xuất đều tuân theo giao diện \texttt{ProposalStrategy} với một REGISTRY trung tâm, qua đó loại bỏ các nhánh \texttt{if/else} lớn từng tồn tại và mở đường cho việc bổ sung loại đề xuất mới chỉ trong một vài tệp riêng lẻ.

\textbf{Về phương pháp luận}, đồ án được thực hiện theo quy trình phát triển phần mềm hiện đại: viết kiểm thử song song với mã nguồn, phát hiện và sửa lỗi rule trước khi đẩy vào nhánh chính, áp dụng các tài liệu pháp lý và nội bộ làm căn cứ cho yêu cầu nghiệp vụ. Các quyết định kỹ thuật quan trọng --- chọn Prisma thay vì Sequelize, áp dụng layered architecture có tầng Repository riêng, dùng Zod làm thư viện kiểm tra dữ liệu thống nhất cho cả backend và frontend --- đều được phân tích trong các chương trước với lý do gắn liền với đặc thù của PM QLKT.

Bên cạnh các kết quả đạt được, đồ án còn tồn tại một số hạn chế. Thứ nhất, phần triển khai mới dừng ở môi trường mạng nội bộ Học viện trên một máy chủ vật lý, chưa đánh giá được khả năng vận hành ở quy mô nhiều cơ sở. Thứ hai, các báo cáo thống kê hiện chủ yếu dạng bảng và biểu đồ tĩnh, chưa có công cụ phân tích so sánh giữa các đơn vị theo nhiều năm liên tiếp. Thứ ba, hệ thống chưa tích hợp với chữ ký số quân đội, do đó các quyết định khen thưởng vẫn cần tải lên dạng PDF đã ký bằng tay.

Quá trình thực hiện đồ án giúp em hiểu sâu hơn về cách trừu tượng hoá quy tắc nghiệp vụ phức tạp thành đối tượng cấu hình thuần tuý, cách thiết kế kiến trúc phân tầng có Repository và lý do tách biệt nó khỏi tầng Service, cách viết kiểm thử bao phủ các kịch bản chu kỳ trượt, cũng như cách áp dụng các văn bản pháp lý và nội bộ vào thiết kế kỹ thuật.

\section{Hướng phát triển}
\label{section:6.2}

Trên nền tảng đã xây dựng, đề xuất bốn hướng phát triển cho phiên bản kế tiếp.

\textbf{(i) Mở rộng các danh hiệu chuỗi cao hơn BKTTCP.} Hiện tại hệ thống dừng ở Bằng khen của Thủ tướng Chính phủ; phiên bản tiếp theo sẽ bổ sung Anh hùng Lực lượng Vũ trang nhân dân, Anh hùng Lao động và các danh hiệu nhà nước cao hơn. Việc mở rộng khả thi nhờ thiết kế cấu hình \texttt{ChainAwardConfig} đã có, chỉ cần bổ sung phần tử mới vào mảng cấu hình kèm hàm xét điều kiện tương ứng.

\textbf{(ii) Bổ sung các loại khen thưởng ngoài bảy nhóm hiện tại.} Ngoài bảy nhóm đã hỗ trợ, thực tế công tác thi đua tại Học viện và các đơn vị quân đội còn nhiều loại khen thưởng khác như khen thưởng theo chuyên đề (an toàn giao thông, học tập và làm theo tư tưởng Hồ Chí Minh), khen thưởng đột xuất theo đợt cao điểm, khen thưởng đối ngoại quốc phòng và các huy chương hữu nghị quốc tế. Kiến trúc Strategy đã chuẩn hoá cho bảy loại hiện tại cho phép mở rộng bằng cách thêm một strategy mới và đăng ký vào REGISTRY trung tâm, không phải sửa lại tầng phê duyệt chung.

\textbf{(iii) Triển khai cho nhiều đơn vị trong toàn quân.} Phiên bản hiện tại được thiết kế cho riêng Học viện Khoa học Quân sự với cây tổ chức hai cấp (CQĐV -- ĐVTT). Hướng tiếp theo là mở rộng mô hình tổ chức thành nhiều cấp (Quân khu / Quân đoàn / Học viện / Trường / Đơn vị cơ sở) và hỗ trợ multi-tenancy để các đơn vị khác trong toàn quân có thể dùng chung hạ tầng phần mềm mà vẫn cách ly dữ liệu nghiệp vụ. Khi đó hệ thống sẽ trở thành nền tảng dùng chung cấp Bộ Quốc phòng thay vì giải pháp riêng của một học viện.

\textbf{(iv) Tích hợp ký số quân đội.} Bổ sung mô-đun ký số PDF quyết định khen thưởng theo chuẩn PKCS\#7 với chữ ký số do hệ thống PKI của Bộ Quốc phòng cấp. Khi đó, file PDF tải lên sẽ chứa chữ ký số hợp lệ, có thể xác thực ngược về cơ quan ký, loại bỏ nhu cầu ký giấy và lưu kho bản giấy.

Bốn hướng được sắp xếp theo độ ưu tiên giảm dần. Hướng (i), (ii) và (iv) khả thi nhất nhờ kiến trúc cấu hình + Strategy đã có và có thể bắt tay trong sáu tháng kế tiếp. Hướng (iii) cần thời gian dài hơn do phụ thuộc vào quyết định triển khai mở rộng phạm vi sử dụng.

\end{document}
